\scanpage{110}

On trouve, en résumé

\[w_*(f_V(\beta_0,\ldots,\beta_{n-1}))=\left(\prod_{0\leq i\leq n-2}(1+\partial(\widetilde\rho_{i,n-1}\widetilde\rho_{i+1,n-1}))\right)\left(1+\partial\bigl(\underbrace{\beta_{n-1}(4-\beta_{n-1})}_{(2+\alpha_{n-1})(2-\alpha_{n-1})}\bigr)\right)\]

formule valable si tous les $\rho_{i,n-1}$ $(0\leq i\leq n-1)$ sont inversibles, ainsi que les $\beta_i$ $(0\leq i\leq n-1)$

Bien sûr, on aurait pu conduire le calcul en sens inverse, en mettant à part $(a_0\ldots a_{n-1})$ etc, \deleted{on} trouve-t-on d'autres expressions remarquables ?

Calculons $w_*(f_{\unsure{E|\mathcal E}})$, \deleted{si on suppose} en notant que $V^\perp=\unsure{\underline O_S|\mathcal O_S}\pi$ et $f(\pi)=-\widetilde\rho\widetilde\sigma$ et on trouve donc un facteur supplémentaire $1+\partial(\underbrace{-\widetilde\rho\widetilde\sigma}_{-\widetilde\rho_{0,n-1}\widetilde\rho_{1,n-1}})$. Posons donc

\[\begin{aligned}\partial(-1)&=\varepsilon&&&&\text{éléments de}\\ \partial(\widetilde\rho_{i,n-1})&=\unsure{\eta_i|y_i}&&0\leq i\leq n-1&& H^1(-,\mathbb Z/2)\\\partial(\widetilde\beta_{n-1})&=\eta_n\end{aligned}\]

On trouve

\[w_*(f_V)\deleted{\unsure{)}}=\prod_{0\leq i\leq n-1}(1+\eta_i+\eta_{i+1})\qquad\text{NB } w_{n+1}(f_V)=0\ !\]

\[w_*(f_E)=(1+\varepsilon+\eta_0+\eta_1)\prod_{0\leq i\leq n-1}(1+\eta_i+\eta_{i+1})\]

% archive sheet number: 109
